% Copertina completa per Amazon KDP: fronte + dorso + retro con smarginatura (bleed).
% Dimensioni trim: 6" x 9" (15,24 x 22,86 cm). Bleed: 0,125" (0,3175 cm) per lato.
% Il dorso dipende dal numero di pagine: viene calcolato da build-copertina-kdp.ps1
% (carta bianca KDP: pagine x 0,002252 pollici) e scritto in spine-kdp.tex.
\documentclass[11pt]{article}

\IfFileExists{spine-kdp.tex}{\def\spinewidthcm{1.5787 cm}
}{\def\spinewidthcm{1.28cm}}

\usepackage{fontspec}
\usepackage{xcolor}
\usepackage{tikz}
\usepackage{geometry}
\usetikzlibrary{arrows.meta,calc}

\definecolor{coverNavy}{HTML}{101B36}
\definecolor{coverBlue}{HTML}{173B68}
\definecolor{coverCyan}{HTML}{55D6D2}
\definecolor{coverOrange}{HTML}{FF9D5C}
\definecolor{coverPale}{HTML}{DCE8F4}
\colorlet{coverLine}{coverCyan}

% Dimensioni della copertina completa.
\newlength{\coverwidth}
\newlength{\coverheight}
\setlength{\coverheight}{23.495cm}
\setlength{\coverwidth}{\dimexpr 0.635cm + 15.24cm + 15.24cm + \spinewidthcm\relax}
\geometry{paperwidth=\coverwidth,paperheight=\coverheight,margin=0pt}

\pagestyle{empty}

\begin{document}

\begin{tikzpicture}[remember picture,overlay]

  % ---- Coordinate (cm): [bleed][retro][dorso][fronte][bleed] ----
  \def\bleed{0.3175cm}
  \def\backw{15.24cm}
  \def\frontw{15.24cm}
  \def\frontH{22.86cm}
  \def\coverH{23.495cm}

  % Sfondo unico su tutta la copertina (si estende fino al bleed).
  \fill[coverNavy] (current page.south west) rectangle (current page.north east);

  % ---- Dorso ----
  \fill[coverBlue]
    ([xshift=\bleed,yshift=\bleed]current page.south west) rectangle
    ([xshift={\bleed+\backw+\spinewidthcm},yshift={\coverH-\bleed}]current page.south west);

  \node[rotate=90,text=white,font=\bfseries\fontsize{16}{20}\selectfont]
    at ([xshift={\bleed+\backw+0.5*\spinewidthcm},yshift={0.5*\coverH}]current page.south west)
    {QUADERNO AI \quad\textbullet\quad 1};

  \node[rotate=90,text=coverPale,font=\small]
    at ([xshift={\bleed+\backw+0.5*\spinewidthcm},yshift={0.5*\coverH - 6.2cm}]current page.south west)
    {Andrea Merlin};

  % ---- Retro ----
  \begin{scope}[shift={([xshift=\bleed,yshift=\bleed]current page.south west)}]

    \node[anchor=north west,text=white,font=\bfseries\fontsize{18}{22}\selectfont]
      at (1.5, 21.3) {Quaderno AI};

    \node[anchor=north west,text=coverPale,font=\footnotesize]
      at (1.5, 19.7) {Concetti, strumenti e percorsi pratici\\per lo sviluppo software con AI};

    \draw[coverOrange,line width=1.4pt] (1.5, 18.3) -- (3.5, 18.3);

    \node[anchor=north west,text=coverPale,font=\small,text width=11.6cm,align=left]
      at (1.5, 17.3)
      {Un percorso operativo per costruire applicazioni .NET con l'intelligenza
       artificiale: server MCP, persistenza con Entity Framework Core, sviluppo
       guidato dalle specifiche, personalizzazione degli agenti e migrazione da
       altri linguaggi.};

    \node[anchor=north west,text=coverCyan,font=\bfseries\fontsize{11}{14}\selectfont]
      at (1.5, 12.6) {Dentro questo quaderno};

    \node[anchor=north west,text=coverPale,font=\footnotesize,text width=11.6cm,align=left]
      at (1.5, 11.6)
      {\textbullet\ Costruire MCP server completi in .NET 10\\
       \textbullet\ Modellare dati e query con EF Core\\
       \textbullet\ Applicare lo spec-driven development con GitHub Spec Kit\\
       \textbullet\ Istruzioni, skill, agenti e plugin per gli agenti di codifica\\
       \textbullet\ Migrare a .NET da PHP, Python, Node.js e Java};

    \draw[coverCyan,opacity=0.55,line width=0.7pt] (1.5, 6.2) -- (7.2, 6.2);

    \node[anchor=north west,text=coverPale,font=\footnotesize,text width=10.4cm,align=left]
      at (1.5, 5.5)
      {Andrea Merlin. Laureato in informatica, appassionato di tecnologie\\
       Microsoft, DevOps e data-privacy; dal 2013 scrive sul suo blog\\
       indipendente di programmazione https://amerlin.keantex.com .\\
       E' speaker e formatore.};

    % Area riservata al codice a barre ISBN (aggiunto da Amazon).
    \fill[white] (9.5, 0.5) rectangle (14.6, 3.55);

  \end{scope}

  % ---- Fronte ----
  \begin{scope}[shift={([xshift={\bleed+\backw+\spinewidthcm},yshift=\bleed]current page.south west)}]

    \fill[coverBlue] (9.00, 0) -- (\frontw, \frontH) -- (\frontw, 9.01) -- cycle;

    \fill[coverNavy] (0, \frontH) -- (15.24, \frontH) -- (15.24, 11.08) -- (0, 8.08) -- cycle;

    \foreach \x/\y/\r in {0.80/1.69/0.58, 2.32/3.62/0.40, 13.35/4.31/0.51,
                          13.93/13.70/0.33, 1.52/18.09/0.44, 12.92/19.55/0.36} {
      \draw[coverLine,opacity=0.28,line width=0.4pt] (\x,\y) circle (\r cm);
    }

    \node[anchor=north west,align=left,text=white]
      at (1.30, 21.40) {%
        {\footnotesize\color{coverCyan}SVILUPPO SOFTWARE \textbullet\ IA \textbullet\ .NET}};

    \draw[coverCyan,opacity=0.65,line width=0.7pt] (1.30, 20.32) -- (4.20, 20.32);

    \coordinate (srcA) at (3.05, 17.47);
    \coordinate (srcB) at (4.43, 16.09);
    \coordinate (srcC) at (3.41, 14.39);
    \coordinate (srcD) at (5.15, 13.24);
    \coordinate (core) at (9.43, 15.32);
    \coordinate (dotnet) at (12.41, 15.32);

    \draw[coverCyan,opacity=0.72,line width=1.1pt]
      (srcA) .. controls +(1.74,0.15) and +(-1.60,1.39) .. (core);
    \draw[coverCyan,opacity=0.55,line width=1.1pt]
      (srcB) .. controls +(1.60,0) and +(-1.45,0.46) .. (core);
    \draw[coverCyan,opacity=0.65,line width=1.1pt]
      (srcC) .. controls +(1.81,0) and +(-1.45,-0.38) .. (core);
    \draw[coverCyan,opacity=0.5,line width=1.1pt]
      (srcD) .. controls +(1.52,-0.15) and +(-1.60,-1.31) .. (core);
    \draw[coverOrange,line width=1.4pt,-{Stealth[length=3mm]}] (core) -- (dotnet);

    \foreach \p/\label in {srcA/MCP,srcB/SPEC,srcC/AGENTI,srcD/.NET} {
      \fill[coverNavy] (\p) circle (0.40cm);
      \draw[coverCyan,line width=1.0pt] (\p) circle (0.40cm);
      \node[text=white,font=\scriptsize\bfseries] at (\p) {\label};
    }

    \fill[coverBlue] (core) circle (0.70cm);
    \draw[coverCyan,line width=1.3pt] (core) circle (0.70cm);
    \draw[coverCyan,opacity=0.32,line width=0.7pt] (core) circle (0.93cm);
    \node[text=white,font=\large\bfseries] at (core) {AI};

    \fill[coverOrange] (dotnet) circle (0.60cm);
    \draw[white,line width=0.6pt,opacity=0.85] (dotnet) circle (0.60cm);
    \node[text=white,font=\footnotesize\bfseries] at (dotnet) {.NET};

    \node[anchor=west,text=coverPale,font=\footnotesize]
      at (1.30, 12.39) {CONCETTI, STRUMENTI E PERCORSI PRATICI};

    \node[anchor=west,align=left,text=white,font=\fontsize{26}{30}\selectfont\bfseries]
      at (1.27, 9.93) {Quaderno AI};

    \node[anchor=west,align=left,text=coverPale,font=\fontsize{11}{14}\selectfont]
      at (1.34, 6.85) {Concetti, strumenti e percorsi pratici\\per lo sviluppo software con AI};

    \draw[coverOrange,line width=1.4pt] (1.34, 4.31) -- (3.48, 4.31);

    \node[anchor=west,text=white,font=\footnotesize]
      at (1.34, 3.54) {MCP \quad SPEC-DRIVEN DEVELOPMENT \quad AGENTI \quad MIGRAZIONE A .NET};

    \node[anchor=south east,text=coverPale,font=\footnotesize]
      at ({\frontw-1.30}, 1.20) {Andrea Merlin \quad | \quad Numero 1 - Novembre 2026};

  \end{scope}

\end{tikzpicture}

\end{document}
